% 三视图：四棱锥，正视底3分三段1、高3，侧视底1，俯视为带交叉对角线的平行四边形；12pt，.6pt灰度
% 点修：三视图改为横向并排（正视图、侧视图、俯视图），图名居中、基线对齐；投影边线/尺寸/虚实未动
\begin{tikzpicture}[line width=.6pt,font=\normalsize,x=1cm,y=1cm,>=Stealth]
% 正视图
\begin{scope}[shift={(0.6,0)}]
\draw (0,0) -- (1.5,3.0) -- (3.0,0) -- cycle;
\draw (1.5,3.0) -- (2.0,0);
\draw[dashed] (1.5,3.0) -- (1.0,0);
% 底边三段 1
\foreach \x in {0,1,2,3} {\draw[thin] (\x,0) -- (\x,-0.45);}
\draw[thin,<->] (0,-0.3) -- (1,-0.3);
\draw[thin,<->] (1,-0.3) -- (2,-0.3);
\draw[thin,<->] (2,-0.3) -- (3,-0.3);
\node[below] at (0.5,-0.35) {1};
\node[below] at (1.5,-0.35) {1};
\node[below] at (2.5,-0.35) {1};
% 左侧高 3（上下引出线分别连到顶点高度与底角）
\draw[thin] (-0.55,3.0) -- (1.5,3.0);
\draw[thin] (-0.55,0) -- (0,0);
\draw[thin,<->] (-0.4,3.0) -- (-0.4,0);
\node[left] at (-0.55,1.5) {3};
\node[below] at (1.5,-0.95) {正视图};
\end{scope}
% 侧视图：窄三角形，底1，高与正视平齐
\begin{scope}[shift={(4.4,0)}]
\draw (0,0) -- (0.5,3.0) -- (1.0,0) -- cycle;
\draw[thin] (0,0) -- (0,-0.45);
\draw[thin] (1.0,0) -- (1.0,-0.45);
\draw[thin,<->] (0,-0.3) -- (1.0,-0.3);
\node[below] at (0.5,-0.35) {1};
\node[below] at (0.5,-0.95) {侧视图};
\end{scope}
% 俯视图：平行四边形及两条交叉对角线
\begin{scope}[shift={(6.6,-0.4)}]
\draw (0,0) -- (3.0,0) -- (4.0,1.2) -- (1.0,1.2) -- cycle;
\draw (0,0) -- (4.0,1.2);
\draw (1.0,1.2) -- (3.0,0);
\node[below] at (2.0,-0.55) {俯视图};
\end{scope}
\end{tikzpicture}
